\clearpage
\fancyhead[L]{\small NOTICES.}
\SourceLeaf{100}
\pdfbookmark[0]{Notices}{anzeigen}
\begin{center}
{\Large\bfseries NOTICES.}\par
\vspace{1.8em}
\rule{.88\textwidth}{.3pt}\par
Göttingische gelehrte Anzeigen.\quad 26 February 1821.\par
\rule{.88\textwidth}{.3pt}
\end{center}
\pdfbookmark[1]{26 February 1821}{anzeige-1821-02-26}

On 15 February, a lecture was submitted to the Royal Society by Councillor \textsc{Gauss}, entitled
\begin{center}
\textit{Theoria Combinationis observationum erroribus minimis obnoxiae, pars prior,}
\end{center}
\noindent and dealing with one of the most important applications of the calculus of probability. No observations concerning the determination of magnitudes in the natural world, however accurately they may be made and however excellent the instruments used, can ever have \textit{absolute} accuracy: they always remain approximations, subject to greater or smaller errors. The errors considered here are not those whose sources are exactly known and whose magnitude can each time be calculated for particular observations; such errors can and should be subtracted from the observed quantities, which is equivalent to their not being present at all. The situation is quite different for errors to be regarded as random, arising from the limited sharpness of the senses, from various unavoidable and irregular imperfections of instruments, and from external disturbances that act irregularly (at least as far as we can tell)—for example, the movement of the atmosphere when looking, or the lack of absolute stability when instruments are set up. These random errors, which
\SourceLeaf{101}
cannot be subjected to calculation, cannot be \textit{removed}; by careful attention and by multiplying the observations, the observer can only \textit{reduce} them. Once the observer has done what is possible, however, it is the geometer's task to assess, by strict mathematical principles, the uncertainty of the observations and of the quantities derived from them by calculation. Most important, wherever quantities connected with the observations can be derived from them by different combinations, the geometer must prescribe the method that leaves the least uncertainty to be feared.

Although random errors as such follow no law, and occur without order—larger in one observation, smaller in another—it is nevertheless certain that, for a specified kind of observation and with the observer and instruments also regarded as fixed, the errors arising from each simple source of error are not merely confined within certain limits: every possible error between those limits also has its own relative probability, so that errors of different sizes are expected more or less often according to their magnitude. Anyone with an exact and complete understanding of the nature of such a source of error would be able to determine these limits and the relation between the probability of each error and its magnitude, in much the same way that, once the rules of a game of chance are known, one can calculate the limits of possible gains and losses and their relative probabilities. The same holds for the total error arising from the combined action of the simple sources of error. These ideas are not confined to direct observations; they also apply to indirect determinations of magnitudes derived from observations. In practice, of course, we will almost always lack the means to specify \textit{a priori} the law governing the probabilities of errors.

How we are to assess, in general, the inadmissibility of a particular kind of observation depends in part on our discretion. We may take only the magnitude of the greatest possible errors as our measure, or also take account of the greater or lesser probability of each possible error. The latter seems more appropriate. This consideration, however, can be brought to bear in many ways
\SourceLeaf{102}
on the matter. Calculators have hitherto used the so-called probable (not \textit{most probable}) error as their measure: it is the error beyond which all possible errors together have just as much probability as all those on this side of it. For this purpose, however, it will be \textit{far more advantageous} to use the \textit{mean} error instead of the probable error, provided that this still-indeterminate notion is defined in the right way. Assign each error a moment depending on its magnitude; multiply the moment of each possible error by its probability and add the products. The error whose moment equals this aggregate must then be regarded as the mean error. Which function of the magnitude of the error we choose for its moment remains, in turn, a matter of our \textit{discretion}, provided that the function is always positive and is greater for larger errors than for smaller ones. The author chose the simplest such function, namely the square; this choice also has several other highly important advantages that no other choice has. Otherwise, any other power with an even exponent could be used; the larger the exponent, the closer one would come to the principle in which only the greatest errors serve as the measure of accuracy. Against the way in which a great geometer defined the mean error—equating the moments of errors to the errors themselves when positive, and using their opposites when negative—it may be observed that this conflicts with mathematical continuity, is no less arbitrary than any other choice, and gives results that are far less simple and satisfactory. It also seems more natural in itself for the moment of errors to grow more strongly than the errors themselves, since one would surely rather tolerate the simple error twice than the double error once.

These explanations had to be given first if anything at all was to be said here about the contents of the investigation, of which the present treatise forms the first part.

When the quantities whose values are found by observation are related in a known way to an equal number of unknown quantities, the values of the unknown quantities\SourceLeaf{103} can, generally speaking, be derived from the observations by calculation. Those values, of course, will be only approximately correct, insofar as the observations were; the calculus of probability has nothing to do here except assess the uncertainty of those determinations, taking the uncertainty of the observations as given. If there are more unknown quantities than observations, the unknowns cannot be determined from the observations at all. But if there are fewer unknown quantities than observations, the problem is overdetermined: infinitely many combinations are then possible for deriving the unknown quantities from the observations. These combinations would all have to yield the same results if the observations were absolutely accurate, but under the circumstances they yield results that differ more or less from one another. Choosing, from this infinitely varied set of combinations, the most suitable one—that is, the one for which the uncertainty of the results is as small as possible—is unquestionably one of the most important problems in applying mathematics to the natural sciences.

The author of the present treatise, who first investigated this problem in 1797 on the principles of the calculus of probability, soon found that it was impossible to determine the \textit{most probable} values of the unknown quantities unless the function representing the probability of errors was known. If it is not known, nothing remains but to assume such a function hypothetically. It seemed most natural to him first to take the reverse course and seek the function that must be assumed if it is to yield, in the simplest of all cases, a rule generally recognized as good: namely, that the arithmetic mean of several values found by observations of equal reliability for one and the same unknown quantity must be regarded as the most probable. It followed that the probability of an error \(x\) must be assumed proportional to an exponential quantity of the form \(e^{-hhxx}\), and that the very method to which he had already been led a few years earlier by other considerations then became generally necessary. He later had occasion to use this method almost daily, especially from 1801 onward, in all kinds of astronomical calculations; \textsc{Legendre} had meanwhile also arrived at it.
\SourceLeaf{104}
This method is now in general use under the name method of least squares; its justification by the calculus of probability, as well as the determination of the accuracy of the results themselves and other related investigations, is developed in detail in the \textit{Theoria Motus Corporum Coelestium}.

The Marquis \textsc{Delaplace}, who later considered this subject from a new point of view—not by seeking the most probable values of the unknown quantities, but by seeking the most suitable combination of the observations—found the remarkable result that, when the number of observations is regarded as infinitely large, the method of least squares is always the most suitable combination, independently of the function expressing the probability of errors.

It can be seen from this that both grounds for the method still leave something to be desired. The first depends entirely on the hypothetical form of the probability of errors; if that form is rejected, the values of the unknown quantities found by the method of least squares are no longer the most probable, just as the arithmetic means in the simplest case mentioned above are no longer so. The second ground leaves us completely in the dark about what is to be done with a moderate number of observations. In that case, the method of least squares no longer has the standing of a law prescribed by the calculus of probability; it recommends itself only through the simplicity of the operations involved.

In the present treatise, the author has taken up this investigation anew, starting from a point of view similar to \textsc{Delaplace}'s but defining the mean error to be feared in a different way, which seems to him more natural in itself. He hopes that friends of mathematics will be pleased to see that, on the new foundation given here, the method of least squares appears in general as the most suitable combination of observations—not approximately, but with mathematical exactness—whatever function may describe the probability of errors and whether the number of observations be large or small.


Many other remarkable investiga\-\SourceLeaf{105}tions are closely connected with the main subject, but their scope compelled the author to reserve the development of most of them for a future second lecture. Of those already included in this first part, allow us to state only one result here. If the function expressing the relative probability of each individual error is unknown, then naturally it is also impossible to determine the probability that an error falls between given limits. Nevertheless, provided that larger errors always have a smaller (at least, not a greater) probability than smaller ones, the probability that the error falls between the limits \(-x\) and \(+x\) must necessarily be greater (at least, not smaller) than \(\frac{x}{m}\surd\frac13\), when \(x\) is less than \(m\surd\frac34\), and not less than \(1-\frac{4mm}{9xx}\), when \(x\) is greater than \(m\surd\frac34\), where \(m\) denotes the mean error to be feared in the observations. For \(x=m\surd\frac34\), the two expressions, as one sees, coincide.
\par\bigskip\centerline{\rule{.88\textwidth}{.3pt}}




